% RQ5 Theoretical Framework: Inter-DAO Cooperation

We develop formal models of multi-DAO dynamics and cooperation mechanisms.

\subsection{Multi-DAO System Model}

\subsubsection{DAO Network}

A DAO network $\mathcal{N} = (\mathcal{D}, \mathcal{R})$ consists of:
\begin{itemize}
    \item $\mathcal{D} = \{d_1, ..., d_n\}$: Set of DAOs
    \item $\mathcal{R} \subseteq \mathcal{D} \times \mathcal{D}$: Cooperation relationships
\end{itemize}

Each DAO $d_i$ has:
\begin{itemize}
    \item Treasury $T_i$
    \item Membership $A_i$
    \item Preferences $\rho_i$ over outcomes
    \item Governance configuration $G_i$
\end{itemize}

\subsubsection{Overlapping Membership}

Agents may hold tokens in multiple DAOs:

\begin{equation}
A_{\text{overlap}} = \{a : a \in A_i \cap A_j \text{ for some } i \neq j\}
\end{equation}

Overlapping membership creates natural preference alignment channels.

\subsection{Joint Proposal Model}

\subsubsection{Inter-DAO Proposals}

An inter-DAO proposal $p^*$ requires approval from a coalition $C \subseteq \mathcal{D}$:

\begin{equation}
\text{pass}(p^*) = \bigwedge_{d \in C} \text{approve}_d(p^*)
\end{equation}

Each DAO votes according to its internal governance rules.

\subsubsection{Resource Contribution}

Joint proposals may require resource contributions:

\begin{equation}
\text{contribution}_d(p^*) = \alpha_d \cdot \text{cost}(p^*)
\end{equation}

where $\alpha_d$ is DAO $d$'s share (negotiated or formula-based).

\subsection{Cooperation Benefit Model}

\subsubsection{Surplus from Cooperation}

Cooperation creates surplus when complementarities exist:

\begin{equation}
S(C) = V(C) - \sum_{d \in C} V(\{d\})
\end{equation}

where $V(C)$ is the value achievable by coalition $C$ and $V(\{d\})$ is the standalone value.

Positive surplus ($S(C) > 0$) motivates cooperation.

\subsubsection{Surplus Division}

The Shapley value provides fair surplus division:

\begin{equation}
\phi_d(V) = \sum_{C \subseteq \mathcal{D} \setminus \{d\}} \frac{|C|!(|\mathcal{D}| - |C| - 1)!}{|\mathcal{D}|!} [V(C \cup \{d\}) - V(C)]
\end{equation}

\subsection{Cooperation Scenarios}

\subsubsection{Homogeneous DAOs}

Similar size, preferences, and capabilities:
\begin{itemize}
    \item Symmetric bargaining power
    \item Natural equal-split norms
    \item Lower coordination cost
    \item But smaller surplus (less complementarity)
\end{itemize}

\subsubsection{Heterogeneous DAOs}

Different sizes, diverse preferences:
\begin{itemize}
    \item Larger potential surplus (complementarity)
    \item Asymmetric bargaining positions
    \item Complex negotiation
    \item Risk of exploitation
\end{itemize}

\subsubsection{Asymmetric Relationships}

Large DAO cooperating with small DAO:
\begin{itemize}
    \item Small DAO vulnerable to domination
    \item Large DAO may free-ride on small DAO effort
    \item Requires explicit fairness mechanisms
\end{itemize}

\subsection{Cooperation Stability}

\subsubsection{Individual Rationality}

Cooperation is individually rational if:

\begin{equation}
\phi_d(V) \geq V(\{d\})
\end{equation}

Each participant gains at least standalone value.

\subsubsection{Coalition Stability}

No subset wants to deviate:

\begin{equation}
\forall S \subset C: \sum_{d \in S} \phi_d(V) \geq V(S)
\end{equation}

\subsection{Theoretical Predictions}

\subsubsection{Prediction 1: Complementarity Drives Surplus}

Cooperation surplus increases with complementarity:

\begin{equation}
S(C) \propto \text{diversity}(C)
\end{equation}

Homogeneous coalitions have low surplus.

\subsubsection{Prediction 2: Size Asymmetry Requires Fairness}

Asymmetric coalitions are unstable without explicit fairness mechanisms.

\subsubsection{Prediction 3: Overlapping Membership Facilitates Cooperation}

Cross-DAO token holders facilitate cooperation by:
\begin{itemize}
    \item Aligning incentives
    \item Transferring information
    \item Building trust
\end{itemize}
